\chapter{Hypothesis 1: Infrastructure and the Access--Usage Gap}
\label{ch:h1}

The hypothesis is tested on three panels, built in the order they appear here. Tracks~1 and~2
measure infrastructure physically, as branches and ATMs; both fail, and
Section~\ref{sec:h1-gate} shows they were always going to. Track~3 substitutes the infrastructure
that actually moved over the sample window --- digital payments --- and recovers the pattern the
hypothesis predicts. The distinction between what was pre-committed and what came afterwards is
load-bearing throughout, and Section~\ref{sec:h1-exploratory} states it plainly rather than
leaving it to a footnote.

\begin{table}[H]
\centering
\caption{The three tracks}
\label{tab:tracks}
\small
\begin{tabular}{@{}llccl@{}}
\toprule
Track & Window & $N$ & Infrastructure measured & Outcome \\
\midrule
1 & \panelA{} & 198 & Branches, ATMs & Fails; cannot be tested \\
2 & \panelB{} & 384 & Branches only & Fails; fragility removed \\
3 & FY2019--FY2023 & 165 & Branches, ATMs, \textbf{UPI} & Predicted pattern recovered \\
\bottomrule
\end{tabular}
\end{table}

\section{Hypothesis and pre-commitment}
\label{sec:h1-hypothesis}

\begin{keybox}{Hypothesis 1}
Greater financial infrastructure raises financial inclusion, but its effect on \textbf{access} is
stronger than its effect on \textbf{usage}.

\medskip
\emph{Rationale.} Opening an account is a supply-side decision that banks can be pushed to make.
Using an account is a demand-side decision shaped by income, trust, and informality.
Infrastructure removes the first constraint, not the second.

\medskip
\emph{Expected:} $\beta_{\text{access}} > \beta_{\text{usage}} > 0$, with the difference
statistically significant.
\end{keybox}

\noindent
The hypothesis is directional twice over: both effects should be \emph{positive}, and the access
effect should be \emph{larger}. Positive and equal effects would reject it while remaining a
substantive result --- genuine deepening rather than formal expansion. Both effects at zero would
be uninformative until the identifying variation had been examined. These readings were tabulated
in advance, and Section~\ref{sec:h1-verdict} returns to that table.

Before any coefficient was estimated, a rule was fixed in writing for what would count as support.

\begin{keybox}{Decision rule, fixed before any result was seen (D10)}
Hypothesis~1 is \textbf{fully supported only if
$\beta_{\text{access}} > \beta_{\text{usage}} > 0$ holds for branches \emph{and} for ATMs}. If it
holds for only one, that is reported as \textbf{partial support}, naming which. This is
pre-committed; it is not to be revised after the coefficients are seen.
\end{keybox}

\noindent
Two features of that rule matter for everything below. It names branches and ATMs specifically ---
UPI is not in it, because UPI was not in the design when it was written. And it was not revised:
the physical verdict reported in Section~\ref{sec:h1-verdict} is the verdict the rule returns.

A second pre-commitment was recorded after the pre-estimation diagnostic failed and before
estimation proceeded: a null result does \textbf{not} license the conclusion ``infrastructure does
not raise access'', being indistinguishable from ``the regressors moved too little to detect an
effect''; any null must be reported with the within-state variation figures attached (D11).

\section{Variables and specification}
\label{sec:h1-variables}

\subsection{Independent variables}

Infrastructure enters in natural units, with no composite index:
\[
\text{Branch}_{it} = \frac{\text{branches}_{it}}{\text{adults}_{it}}\times 10^5,
\qquad
\text{ATM}_{it} = \frac{\text{ATMs}_{it}}{\text{adults}_{it}}\times 10^5,
\qquad
\text{UPI}_{it} = \frac{\text{registered users}_{it}}{\text{adults}_{it}}.
\]

An earlier design combined branches, ATMs and PMJDY accounts into one normalised index. It was
abandoned for three reasons (D3, D10). \textbf{PMJDY had to go, and its removal proved
load-bearing:} PMJDY accounts are an \emph{account count}, so placing them among the regressors
while deposit and credit account counts sit among the outcomes would put an account count on both
sides of the equation, manufacturing exactly the mechanical correlation the design warns against.
\textbf{Separating branches from ATMs is statistically safe:} their raw correlation is 0.904, but
that is almost entirely cross-sectional, and after state and year demeaning --- the variation the
model actually uses --- it falls to 0.305, with a variance inflation factor of 1.10.
\textbf{Separating them makes a prediction an index cannot:} a branch is where an account is
\emph{opened}, an ATM where it is \emph{used}. Abandoning the index also removed two arbitrary
choices, the 0--1 rescaling of the regressors and the branches-versus-ATMs weighting, and left
coefficients in interpretable units.

For Track~3, the digital measure is \textbf{registered UPI users per adult}, in levels. Three
candidates existed. Transactions per capita and value per capita were rejected as regressors
because they are themselves \emph{usage} of the digital channel, and using a flow to explain a flow
of the same kind would be close to circular. Registered users is an adoption stock --- whether the
channel is available and taken up at all --- which is the role branches and ATMs play on the
physical side. Levels rather than logs keeps all three infrastructure regressors in the same kind
of unit (D16).

\subsection{Dependent variables}

Four indicators form two constructs, with denominators the design distinguishes deliberately:
\textbf{access is measured per adult, usage per capita}.

\begin{table}[H]
\centering
\caption{Construction of the dependent variables}
\label{tab:dvconstruction}
\small
\begin{tabular}{@{}lllc@{}}
\toprule
Construct & Indicator & Denominator & Entangled with UPI? \\
\midrule
\multirow{2}{*}{\Access{}}
 & Deposit accounts per adult & Population 18+ & \textbf{Yes} (see \S\ref{sec:h1-decomp}) \\
 & Credit accounts per adult  & Population 18+ & No \\
\addlinespace
\multirow{2}{*}{\Usage{}}
 & Deposits per capita (\rs) & Total population & No \\
 & Credit per capita (\rs)   & Total population & No \\
\bottomrule
\end{tabular}
\end{table}

\noindent
Each indicator is rescaled by pooled min--max normalisation,
$\tilde{x}_{it} = (x_{it} - \min x)/(\max x - \min x)$, with the extremes taken over all
state-years rather than within each year, so that genuine growth over time --- which is what the
fixed-effects model estimates from --- is preserved rather than erased. The constructs are the
simple average of their two components, and the gap is their difference:
$\Access{} = \frac{1}{2}(\tilde{d} + \tilde{c})$,
$\Usage{} = \frac{1}{2}(\tilde{D} + \tilde{C})$,
$\AUG{} = \Access{} - \Usage{}$. The normalisation for Track~3 is held at the Track~1 scaling, so
the dependent variables are identical across tracks and the results stay comparable (D15).

\subsection{Specification}
\label{sec:h1-spec}

Three equations are estimated by two-way fixed effects with standard errors clustered by state.
Writing $\mathbf{X}_{it}$ for the infrastructure regressors --- branches and ATMs in Tracks~1
and~3, branches alone in Track~2, with UPI added in Track~3:

\begin{align}
\Access{}_{it} &= \alpha_i + \lambda_t + \boldsymbol{\beta}'\mathbf{X}_{it}
                  + \gamma \ln\mathrm{NSDP}_{it} + \varepsilon_{it} \label{eq:access}\\[1pt]
\Usage{}_{it}  &= \alpha_i + \lambda_t + \boldsymbol{\delta}'\mathbf{X}_{it}
                  + \gamma \ln\mathrm{NSDP}_{it} + \varepsilon_{it} \label{eq:usage}\\[1pt]
\AUG{}_{it}    &= \alpha_i + \lambda_t + \boldsymbol{\theta}'\mathbf{X}_{it}
                  + \gamma \ln\mathrm{NSDP}_{it} + \varepsilon_{it} \label{eq:aug}
\end{align}

State effects $\alpha_i$ absorb everything fixed within a state --- geography, banking legacy,
2011 baselines --- and year effects $\lambda_t$ absorb national shocks common to all states, which
over this window include COVID-19, PMJDY account opening, and the national growth of UPI. Estimates
are therefore identified from \textbf{within-state variation over time, relative to the national
trend in that year}; all cross-sectional variation is discarded. That is the right specification
for the causal question and, as Section~\ref{sec:h1-gate} shows, the source of the study's central
difficulty. With 33 clusters, cluster-robust standard errors sit at the lower end of the range in
which the asymptotic approximation is reliable \citep{cameron2015}.

\paragraph{Equation~\eqref{eq:aug} is not a third finding.} Since $\AUG{}$ is defined as
$\Access{} - \Usage{}$ and all three equations share a right-hand side, linearity gives
$\theta_k = \beta_k - \delta_k$ exactly. It contains no information beyond the first two equations;
its purpose is to supply the \textbf{standard error on the difference}, which is the formal test of
the hypothesis. The claim is not merely that $\beta > \delta$ in point estimates, but that the gap
is distinguishable from zero. The identity is asserted numerically at every estimation.

\section{The diagnostic that governs everything}
\label{sec:h1-gate}

The design specifies this diagnostic first, and specifies it as a \emph{gate} rather than as
description: if infrastructure barely moves within states, fixed effects have nothing to estimate
from. Thresholds were fixed \textbf{before the numbers were computed} --- a within-state share of
variation below 10 per cent is a FAIL, below 25 per cent WEAK, at or above 25 per cent a PASS.

\begin{table}[H]
\centering
\caption{Variance decomposition. Thresholds fixed before computation:
FAIL $<10\%$, WEAK $<25\%$, PASS $\geq 25\%$}
\label{tab:within}
\footnotesize
\begin{tabular}{@{}lrrrrl@{}}
\toprule
Variable & Overall SD & Between SD & Within SD & Within share & Verdict \\
\midrule
\multicolumn{6}{@{}l}{\textit{Physical regressors, \panelA{}}}\\
\textbf{Branches per lakh adults} & 11.455 & 11.574 & 0.805 & \textbf{7.0\%} & \textbf{FAIL} \\
\textbf{ATMs per lakh adults}     & 18.925 & 19.097 & 1.651 & \textbf{8.7\%} & \textbf{FAIL} \\
\addlinespace
\multicolumn{6}{@{}l}{\textit{Outcomes and control, same panel}}\\
\Access{}              & 0.172 & 0.167 & 0.049 & 28.4\% & --- \\
\Usage{}               & 0.200 & 0.201 & 0.029 & 14.4\% & --- \\
\AUG{}                 & 0.128 & 0.127 & 0.029 & 22.9\% & --- \\
$\ln$(NSDP per capita) & 0.541 & 0.528 & 0.147 & 27.1\% & --- \\
\addlinespace
\multicolumn{6}{@{}l}{\textit{Digital regressors, FY2019--FY2023 (Track 3)}}\\
UPI users per adult          & --- & --- & --- & \textbf{67.1\%} & \textbf{PASS} \\
UPI transactions per capita  & --- & --- & --- & \textbf{79.9\%} & \textbf{PASS} \\
UPI value per capita         & --- & --- & --- & \textbf{78.8\%} & \textbf{PASS} \\
\bottomrule
\end{tabular}
\end{table}

\begin{keybox}{The gate failed for physical infrastructure --- and only for physical infrastructure}
Roughly \textbf{93 per cent} of the variation in branches and ATMs is \emph{between} states, which
is exactly what two-way fixed effects discard. The digital series clear the PASS threshold by a
factor of three to eight.
\end{keybox}

\begin{figure}[H]
\centering
\includegraphics[width=\textwidth]{fig1_variables_over_time.png}
\caption{The identification problem. Each grey line traces one state across \panelA{}; the red line
is the cross-state median. \textbf{Top row, the regressors:} state lines run flat, with median
within-state swings of 5.6 per cent (branches) and 11.8 per cent (ATMs). \textbf{Bottom row, the
outcomes:} the same states climb steadily, swinging 47.9 and 61.9 per cent. Fixed effects use only
the vertical movement of individual lines, and there is very little of it in the top row.}
\label{fig:overtime}
\end{figure}

\subsection{Why the physical variation is so thin}

Three mechanisms compound, diagnosed before any coefficient was examined. \textbf{Infrastructure
stocks are slow-moving:} raw branch counts swing a median of only 10.7 per cent within a state
across six years --- 3.6 per cent in Chandigarh, 3.7 in Punjab. \textbf{The denominator cancels
much of what remains:} adult population grows a median 10.2 per cent within states over the window,
in the \emph{same direction} as branches, so branches per lakh adults swing only 5.6 per cent;
because the adult series is interpolated (Section~\ref{sec:data-defects}) it is almost perfectly
smooth and contributes no independent variation, merely deflating a trend. \textbf{The outcomes
move five to ten times more than the regressors:} credit accounts have a 30.0 per cent within share
against 7.0 for branches. Wide standard errors are guaranteed by construction, whether or not an
effect exists.

The problem is confined to the regressors --- the outcomes and the control have ample within-state
variation. The panel is not too short in general; it is too short \emph{for physical
infrastructure}. Lengthening the window helps but does not cure: the within share for branches
rises from 7.0 per cent over six years to 13.4 over ten, 20.1 over twelve, and 23.2 over thirteen,
still short of the PASS line. ATMs cannot follow at all, since the state-wise release begins in
2018. This is what motivates Track~2, and ultimately what motivates Track~3.

\section{Tracks 1 and 2: physical infrastructure}
\label{sec:h1-physical}

\begin{table}[H]
\centering
\caption{Two-way fixed-effects estimates, physical infrastructure. State and year effects
throughout; standard errors clustered by state. Track~2 drops ATMs, which cannot be extended
before 2018.}
\label{tab:physical}
\footnotesize
\begin{threeparttable}
\begin{tabular}{@{}llrrrrrr@{}}
\toprule
& & \multicolumn{3}{c}{Track 1 --- \panelA{}, $N=198$}
  & \multicolumn{3}{c}{Track 2 --- \panelB{}, $N=384$} \\
\cmidrule(lr){3-5}\cmidrule(lr){6-8}
Equation & Regressor & Coef. & SE & $p$ & Coef. & SE & $p$ \\
\midrule
\multirow{2}{*}{(1) \Access{}}
 & Branches & $-0.00897$ & 0.00699 & 0.201 & $-0.00483$ & 0.00645 & 0.455 \\
 & ATMs     & $-0.00164$ & 0.00377 & 0.664 & \multicolumn{3}{c}{---} \\
\addlinespace
\multirow{2}{*}{(2) \Usage{}}
 & Branches & $\mathbf{-0.01264}$ & 0.00363 & $\mathbf{0.001}$
            & $\mathbf{-0.01032}$ & 0.00570 & $\mathbf{0.071}$ \\
 & ATMs     & $-0.00109$ & 0.00263 & 0.678 & \multicolumn{3}{c}{---} \\
\addlinespace
\multirow{2}{*}{(3) \AUG{}}
 & Branches & $+0.00367$ & 0.00558 & 0.512 & $+0.00549$ & 0.00385 & 0.155 \\
 & ATMs     & $-0.00055$ & 0.00159 & 0.731 & \multicolumn{3}{c}{---} \\
\midrule
\multicolumn{8}{@{}l}{Within-$R^2$ --- Track 1: 0.069, 0.042, 0.111.
 Track 2: $-0.172$, $-0.478$, 0.215.} \\
\multicolumn{8}{@{}l}{Joint test (Track 1, Branch $=$ ATM $=$ 0): \Access{} $p = 0.374$;
 \Usage{} $p = 0.0005$; \AUG{} $p = 0.710$.} \\
\bottomrule
\end{tabular}
\begin{tablenotes}\scriptsize
\item The identity $\theta = \beta - \delta$ holds exactly throughout: for Track~1 branches,
$-0.00897 - (-0.01264) = +0.00367$. Track~2's negative within-$R^2$ values are reported rather than
suppressed; with two-way fixed effects the measure can go negative, and here it indicates the
regressors add nothing over the fixed effects alone.
\end{tablenotes}
\end{threeparttable}
\end{table}

\noindent
Read before interpretation: \textbf{nothing is significant in the \Access{} equation} on either
track, and every coefficient there is \emph{negative}. \textbf{Exactly one coefficient is
significant --- branches on \Usage{}, and it is negative.} \textbf{\AUG{} is not significant} for
any regressor on either track, so the formal test of the hypothesis returns nothing.
\textbf{ATMs are insignificant everywhere.} Explanatory power is very low throughout, exactly as
the gate forecast.

\begin{figure}[H]
\centering
\includegraphics[width=\textwidth]{fig2_coefficients.png}
\caption{\textbf{Left:} every physical-infrastructure coefficient across both tracks, with 95 per
cent confidence intervals. Exactly one --- branches on \Usage{} in Track~1 --- is significant at
the 5 per cent level, and it is negative. \textbf{Right:} the coefficient on branches in the
\AUG{} equation across all four physical specifications. It is positive in every one, in the
direction the hypothesis predicts, but never attains conventional significance: $p$ runs $0.512$,
$0.098$, $0.155$, $0.084$. Reporting this as a finding would mean loosening the pre-committed rule
after seeing the coefficients, so it is recorded as a suggestive pattern only.}
\label{fig:coefs}
\end{figure}

\subsection{Random effects is not an escape}

Random effects would use both within- and between-state variation, and so appears to solve the
problem. The Hausman test asks whether doing so is consistent, under the null that state effects
are uncorrelated with the regressors \citep{hausman1978}. On Track~1 the null is rejected
decisively in all three equations ($\chi^2 = 108.3$, $73.2$ and $12.4$; $p < 0.0001$, $<0.0001$ and
$0.006$). State effects \emph{are} correlated with branches and ATMs --- the ``wealthier states
attract more infrastructure'' channel, confirmed rather than assumed. On Track~2 the evidence is
weaker: \Access{} now favours random effects ($p = 0.164$), while \Usage{} and \AUG{} still favour
fixed effects.

\begin{keybox}{The bind}
The Hausman test says fixed effects are \textbf{required}. The variance decomposition says fixed
effects have almost \textbf{nothing to work with}. Dropping the state fixed effects to recover
precision is not available --- the data itself rejects that specification. The design cannot be
rescued by changing estimator.
\end{keybox}

\subsection{Two observations decide the ATM result}
\label{sec:h1-robustness}

All equations were re-estimated dropping the two known-defective observations identified in
Section~\ref{sec:data-defects} --- Delhi 2023 and Chandigarh 2023, two cells out of 198.

\begin{table}[H]
\centering
\caption{Track~1 estimates against estimates excluding Delhi 2023 and Chandigarh 2023}
\label{tab:robustness}
\footnotesize
\begin{tabular}{@{}llrrrrl@{}}
\toprule
& & \multicolumn{2}{c}{Main ($N=198$)} & \multicolumn{2}{c}{Trimmed ($N=196$)} & \\
\cmidrule(lr){3-4}\cmidrule(lr){5-6}
Equation & Regressor & Coef. & $p$ & Coef. & $p$ & Change \\
\midrule
\Access{} & Branches & $-0.00897$ & 0.201 & $-0.00570$ & 0.078 & \\
\Access{} & \textbf{ATMs} & $\mathbf{-0.00164}$ & 0.664 & $\mathbf{+0.00331}$ & $\mathbf{0.004}$
  & \textbf{sign flip; becomes significant} \\
\addlinespace
\Usage{}  & Branches & $-0.01264$ & \textbf{0.001} & $-0.01185$ & \textbf{0.000} & stable \\
\Usage{}  & \textbf{ATMs} & $-0.00109$ & 0.678 & $+0.00011$ & 0.957 & \textbf{sign flip} \\
\addlinespace
\AUG{}    & Branches & $+0.00367$ & 0.512 & $+0.00615$ & 0.098 & \\
\AUG{}    & \textbf{ATMs} & $-0.00055$ & 0.731 & $+0.00319$ & 0.148 & \textbf{sign flip} \\
\bottomrule
\end{tabular}
\end{table}

\begin{keybox}{Two observations out of 198 reverse the sign of every ATM coefficient}
They also move the estimated effect of ATMs on \Access{} from null ($p = 0.664$) to significant and
positive ($p = 0.004$). A result that depends on two state-years is not a result: \textbf{nothing
about ATMs can be concluded in either direction.}
\end{keybox}

\noindent
This is what the gate predicted. When a regressor carries only 7 to 9 per cent within-state
variation, individual observations acquire enormous leverage. The two dropped cells are moreover
\emph{known defects in RBI's published source}, not outliers chosen for being inconvenient --- so
the headline ATM result had been riding on data already established as wrong. On Track~2, with the
window doubled and within-state variation at 15.3 per cent, the same trim produces \textbf{no sign
flips anywhere}: the fragility disappears with more data, though the verdict does not change.

\section{Track 3: substituting the infrastructure that moved}
\label{sec:h1-digital}

\subsection{Why}

Over \panelA{}, physical banking infrastructure in India was flat while retail payments moved onto
phones. If financial infrastructure was doing anything over this window, it was not doing it
through branches. Figure~\ref{fig:premise} states that premise in numbers.

\begin{figure}[H]
\centering
\includegraphics[width=\textwidth]{fig4_digital_vs_physical.png}
\caption{\textbf{Left:} unweighted state means indexed to FY2019, on a log scale. UPI value and
transactions per capita rise roughly twenty-five-fold in five years; branches and ATMs per lakh
adults are flat. \textbf{Right:} the same in levels --- UPI users per adult rise from 0.135 to
0.464 while branches per lakh adults sit between 21.4 and 22.0 throughout.}
\label{fig:premise}
\end{figure}

\begin{figure}[H]
\centering
\includegraphics[width=0.92\textwidth]{fig5_within_variation.png}
\caption{Within-state share of total variation, against thresholds fixed before the numbers were
computed. The reason Tracks~1 and~2 could not test the hypothesis does not apply to Track~3.}
\label{fig:variation}
\end{figure}

\noindent
The consequence for estimation is immediate. Within-state variation in UPI users per adult is
\textbf{67.1 per cent}, against 6.5 for branches and 8.2 for ATMs on the same sample. Explanatory
power rises correspondingly: within-$R^2$ goes from 0.069 to \textbf{0.609} for \Access{} and from
0.042 to \textbf{0.863} for \Usage{}. Fixed effects plus branches and ATMs explained essentially
nothing; one digital variable explains most of it.

\subsection{The test that mattered}
\label{sec:h1-decomp}

One objection has to be dealt with before any digital coefficient can be read. PhonePe registration
requires a KYC-linked bank account, so UPI adoption and \emph{deposit-account counts} are entangled
by accounting rather than by economics --- the same trap that dropping PMJDY was meant to avoid,
and deposit accounts are half of the \Access{} construct. Credit accounts are not entangled the
same way: a bank loan account is neither required for nor created by PhonePe registration.

The outcome was therefore decomposed into its four underlying indicators, with the reading fixed
before the regressions were run (D17): if UPI $\to$ credit accounts is large and significant, the
\Access{} result is not merely definitional; if credit is null while deposits are large, the
\Access{} coefficient \emph{is} the accounting relationship and \Access{} must be reported
credit-only thereafter; if both are large, only the credit result gets a substantive reading.

\begin{figure}[H]
\centering
\includegraphics[width=\textwidth]{fig3_digital.png}
\caption{\textbf{Left:} the coefficient on UPI users per adult for each of the four outcome
components. The entangled channel is the null one, and collapses further when the two defective
cells are dropped. \textbf{Right:} the hypothesis pattern on the credit-only specification.}
\label{fig:digital}
\end{figure}

\begin{table}[H]
\centering
\caption{Component decomposition: coefficient on UPI users per adult, $N=165$, and the same
estimate excluding Delhi 2023 and Chandigarh 2023}
\label{tab:components}
\footnotesize
\begin{tabular}{@{}llrrrr@{}}
\toprule
Component & Status & Coef. & $p$ & Coef. (trimmed) & $p$ \\
\midrule
Deposit accounts per adult & \textbf{entangled} & $+0.193$ & 0.206 & $+0.018$ & 0.509 \\
Credit accounts per adult  & clean & $\mathbf{+0.439}$ & $\mathbf{0.0001}$ & $+0.445$ & 0.0003 \\
Deposits \rs per capita    & clean & $+0.222$ & $<0.001$ & --- & --- \\
Credit \rs per capita      & clean & $+0.145$ & $<0.001$ & --- & --- \\
\addlinespace
\Access{} (composite)      & mixed & $+0.316$ & 0.0004 & $+0.231$ & 0.0005 \\
\bottomrule
\end{tabular}
\end{table}

\begin{keybox}{The concern was not merely unfounded --- it pointed the wrong way}
The \textbf{entangled} channel is the null one: UPI $\to$ deposit accounts is $+0.193$
($p = 0.206$), collapsing to $+0.018$ ($p = 0.509$) when the two known-defective cells are dropped,
with the standard error falling from 0.151 to 0.027. The \textbf{clean} channel carries everything:
UPI $\to$ credit accounts is $+0.439$ ($p = 0.0001$) and is completely unmoved by the same trim.

\medskip
The entangled component was \emph{diluting} a real credit effect, not manufacturing one. The
composite \Access{} number averages a genuine effect with a defect-driven artefact, so
\textbf{\Access{} is reported credit-only from here}.
\end{keybox}

\subsection{The hypothesis pattern, on the credit-only specification}

\begin{table}[H]
\centering
\caption{Track~3, credit-only \Access{}. Two-way fixed effects, $N = 165$, standard errors
clustered by state}
\label{tab:digital}
\small
\begin{tabular}{@{}lrrr@{}}
\toprule
Regressor & \Access{} (credit accounts) & \Usage{} & Gap (\AUG{}) \\
\midrule
\textbf{UPI users per adult} & $\mathbf{+0.439}$ ($p<0.001$) & $\mathbf{+0.183}$ ($p<0.001$)
  & $\mathbf{+0.256}$ ($\mathbf{p=0.020}$) \\
Branches per lakh adults & $+0.0091$ ($p=0.085$) & $-0.0062$ ($p=0.020$) & $+0.0153$ ($p=0.018$) \\
ATMs per lakh adults & $-0.0007$ ($p=0.773$) & $-0.0002$ ($p=0.868$) & $-0.0005$ ($p=0.884$) \\
\bottomrule
\end{tabular}
\end{table}

\noindent
Four things follow. \textbf{For UPI, $\beta_{\text{access}} > \beta_{\text{usage}} > 0$ and the gap
is significant at 5 per cent} --- the pattern the hypothesis predicts, and the first time in this
study that the access--usage gap is distinguishable from zero. \textbf{Branches turn positive on
access for the first time}, once \Access{} is credit-only and digital adoption is controlled for,
though only at $p = 0.085$. \textbf{ATMs are null on everything}, and the ATM coefficient that
looked significant under trimming in Track~1 is revealed as a composite-outcome artefact.
\textbf{Physical infrastructure still fails the pre-committed rule}, and that verdict never
changed.

\subsection{Robustness}

The headline coefficient was refit 33 times, dropping each state in turn.

\begin{figure}[H]
\centering
\includegraphics[width=0.95\textwidth]{fig6_leave_one_out.png}
\caption{Leave-one-state-out estimates of the coefficient on UPI users per adult in the credit
accounts equation. The estimate stays between $+0.356$ and $+0.477$, and no refit loses
significance at the 5 per cent level. Maharashtra is the most influential single state, and the
result is \emph{stronger} without it.}
\label{fig:loo}
\end{figure}

\noindent
The contrast with Track~1 is the point: there, two cells out of 198 flipped every ATM sign; here,
removing an entire state moves the coefficient by at most 19 per cent and never changes the
conclusion. The result is also unmoved by the D7 data defects, which affect savings accounts and
therefore cannot touch a credit-only outcome.

\subsection{A mechanism that fits}

Digital payment history feeding credit underwriting --- cash-flow-based lending, and the account
aggregator framework --- is a documented development in India over exactly this window, and the
wider literature finds digital payment adoption spreading through network effects rather than
infrastructure stock \citep{crouzet2023,dsilva2019}. It predicts precisely what is observed: UPI
moving loan accounts while leaving deposit accounts alone. That the data is consistent with a
mechanism is suggestive of it, not proof.

\section{Verdict}
\label{sec:h1-verdict}

\subsection{Against the pre-committed rule}

The rule requires $\beta_{\text{access}} > \beta_{\text{usage}} > 0$ for branches \emph{and} ATMs.
For branches on Track~1, the ordering holds ($-0.00897 > -0.01264$) but both coefficients are
negative, so the positivity condition fails; the difference is insignificant ($p = 0.512$). For
ATMs the ordering does not hold, nothing is significant, and every sign reverses on a two-cell
trim. Track~2 reproduces the branch result with twice the power and none of the fragility.

\begin{keybox}{Verdict on the pre-committed hypothesis}
\textbf{Hypothesis 1 is not supported for physical infrastructure.} This is not partial support:
the positivity condition fails for branches, and the ordering condition fails for ATMs. The same
conclusion is reached on all three tracks, and it never changed.
\end{keybox}

\noindent
Mapping onto the outcomes tabulated in advance, the row that applies is \emph{both effects
$\approx 0$ --- check within-state variation before concluding anything}. The diagnostic was run
first, as instructed, and returned 7.0 and 8.7 per cent. The correct conclusion is therefore
\textbf{not} that infrastructure fails to raise financial inclusion, but that \textbf{Tracks~1
and~2 cannot test the question}: the regressors barely move, the estimator the Hausman test
requires discards the variation that does exist, and the estimates are unstable to two
observations. The null is uninformative, exactly as fixed in advance that it would be.

\subsection{What the digital result is, and what it is not}
\label{sec:h1-exploratory}

\begin{keybox}{The digital result is exploratory by construction}
The pre-committed rule named \textbf{branches and ATMs}, and it failed. \textbf{UPI was added after
seeing that null.} The digital finding therefore cannot be reported as a hypothesis that was tested
and passed --- it is a hypothesis \textbf{generated} by this data, for a future design to test.
Stating it any other way is exactly the retro-fitting that pre-committing was meant to prevent.
\end{keybox}

\noindent
This is the first thing a careful reader will test, so it is stated here rather than buried. What
can be said is bounded but not trivial: on a measure of financial infrastructure that genuinely
moved within states, the access--usage gap widens significantly, the effect runs entirely through
the channel that is not definitionally entangled with the regressor, and the result survives the
removal of any single state. The identification problem that defeated Tracks~1 and~2 is solved. The
causal problem is not --- fixed effects remove fixed state differences and the common national
trend, but not state-specific time-varying confounders, and smartphone penetration, fintech entry
and youth share all rose alongside both UPI adoption and credit accounts over these five years.

\subsection{The persistent puzzle}

One physical coefficient is significant and robust across specifications: \textbf{branches on
\Usage{}, negative} ($p = 0.001$ on Track~1, $0.071$ on Track~2, $0.020$ on Track~3). Within a
state, years with more branches per lakh adults show lower \Usage{}. It is the opposite sign to
anything the hypothesis anticipated, and it is treated as a puzzle rather than a result: the
variation problem applies to it as much as to the nulls, \Usage{} is severely compressed by the
pooled normalisation, and a negative within-state relationship is at least as consistent with
branches being \emph{placed} where deposit growth is weak as with branches causing it. Nothing here
identifies direction.

\section{Limitations}
\label{sec:h1-limitations}

Ordered by how much damage each does to the conclusions.

\paragraph{1. Identification --- the binding limitation.} On the physical tracks the regressors
barely move: 7.0 and 8.7 per cent within-state variation, below the FAIL threshold fixed before
looking. The null on branches and ATMs \textbf{cannot be distinguished from ``not enough variation
to detect anything''}, and any such null must be reported with these numbers attached. On the
digital track variation is ample and that objection does not apply --- but fixed effects do not
remove state-specific time-varying confounders, so the digital result is a within-state correlation
between two co-trending series and \textbf{no causal claim is available}. Reverse causality is live
rather than hypothetical: the Hausman test rejecting random effects is direct evidence that state
effects are correlated with the regressors.

\paragraph{2. The digital result is exploratory by construction.} UPI was added after seeing the
physical null (Section~\ref{sec:h1-exploratory}).

\paragraph{3. Known data defects.} The deposit-account breaks in RBI's source (Delhi 2023,
Chandigarh 2023 inside the window) were enough to move the deposit-account coefficient from
$+0.193$ to $+0.018$. Adult population is interpolated rather than published for 243 of 266
state-years, and being near-linear it \emph{suppresses} within-state variation in the per-adult
regressors, making the first limitation worse than it looks. Two units are dropped for want of
income data, and boundary changes forced merges, so nothing state-specific can be said about
either half of a merged unit.

\paragraph{4. Measurement.} \textbf{Attribution differs three ways:} bank deposits and credit are
recorded at branch location, loan accounts by place of sanction, and PhonePe by the user's
registered state --- three different definitions of ``where'', merged into one panel. The first
inflates Maharashtra and Delhi, where lending is booked. \textbf{PhonePe is one firm}, roughly
45--50 per cent of UPI, with market share varying by state, so cross-state levels are noisier than
within-state growth; NPCI's total UPI series would fix this. \textbf{\Usage{} is severely
right-skewed after normalisation} --- 85 per cent of observations below 0.20 --- so
$\beta_{\text{access}}$ and $\beta_{\text{usage}}$ were never on fully equivalent scales, a direct
threat to the comparison itself, left unfixed deliberately in order to test the specification as
written. And the study is \textbf{state-level, not person-level}: nothing here can show that
\emph{individuals} hold unused accounts, which is the mechanism the hypothesis describes.

\paragraph{5. Scope.} The physical panel stops at FY2023, FY2024 having been traded away for a
consistent bank universe; the digital panel is five years, since PhonePe data begin in 2018; ATMs
cannot extend before 2018; and dropping PMJDY, though necessary, left the physical measure with
nothing digital in it until PhonePe was added.

\begin{keybox}{What survives all of this}
The \textbf{physical null}, reported as uninformative rather than as evidence of no effect.
\textbf{Branches $\to$ \Usage{} negative} ($p \leq 0.020$ across three specifications and two
samples), robust but with no causal reading available. \textbf{UPI $\to$ credit accounts positive}
($+0.439$), robust to dropping any single state and to the known data defects --- but
correlational. And \textbf{UPI $\to$ deposit accounts $\approx 0$}, which is the asymmetry the whole
digital finding rests on.
\end{keybox}
